\section{Introduction}
\label{sec:intro}
Sparsity is a recurring design principle for intelligent systems, spanning from biological neural circuits~\citep{olshausen1997sparse,lennie2003cost} to modern Large Language Models (LLMs). Currently, this principle is primarily realized through Mixture-of-Experts (MoE)~\citep{shazeer2017outrageously,dai2024deepseekmoe}, which scales capacity via conditional computation. 
Owing to its ability to drastically increase model size without proportional increases in compute, MoE has become the de facto standard for frontier models~\citep{guo2025deepseek,comanici2025gemini,kimi-k2}.


Despite the success of this conditional computation paradigm, the intrinsic heterogeneity of linguistic signals suggests significant room for \textit{structural optimization}. Specifically, language modeling entails two qualitatively different sub-tasks: compositional reasoning and knowledge retrieval. While the former demands deep, dynamic computation, a substantial portion of text—such as named entities and formulaic patterns—is local, static, and highly stereotyped~\citep{erman2000idiom,constant2017survey}. The effectiveness of classical $N$-gram models~\citep{liu2024infini,nguyen2024understanding,brants-etal-2007-large} in capturing such local dependencies implies that these regularities are naturally represented as computationally inexpensive lookups. Since standard Transformers~\citep{vaswani2017attention} lack a native knowledge lookup primitive, current LLMs are forced to \textbf{simulate retrieval through computation}. For instance, resolving a common multi-token entity requires consuming multiple early layers of attention and feed-forward networks~\citep{ghandeharioun2024patchscopes,DBLP:conf/coling/JinYHZWH0MMDYDZ25} (see \autoref{tab:patchscope}). This process essentially amounts to an expensive runtime reconstruction of a static lookup table, wasting valuable sequential depth on trivial operations that could otherwise be allocated to higher-level reasoning.

To align model architecture with this linguistic duality, we advocate for a complementary axis of sparsity: conditional memory. Whereas conditional computation sparsely activates parameters to process dynamic logic~\citep{shazeer2017outrageously,bengio2013estimatingpropagatinggradientsstochastic}, conditional memory relies on sparse lookup operations to retrieve static embeddings for fixed knowledge. As a preliminary exploration of this paradigm, we revisit $N$-gram embeddings~\citep{bojanowski2017enriching} as a canonical instantiation: local context serves as a key to index a massive embedding table via constant-time $\mathcal{O}(1)$ lookups~\citep{tito2017hash,huang2025over,pagnoni2025byte,yu2025scaling}. Our investigation reveals that, perhaps surprisingly, this static retrieval mechanism can serve as an ideal complement to modern MoE architecture—but only if it is properly designed. In this paper, we propose \textbf{Engram}, a conditional memory module grounded in the classic $N$-gram structure but equipped with modern adaptations such as tokenizer compression, multi-head hashing, contextualized gating, and multi-branch integration~(detailed in \autoref{sec:architecture}).

To quantify the synergy between these two primitives, we formulate the \textit{Sparsity Allocation} problem: given a fixed total parameter budget, how should capacity be distributed between MoE experts and Engram memory? Our experiments uncover a distinct U-shaped scaling law, revealing that even simple lookup mechanisms, when treated as a first-class modeling primitive, act as essential complements to neural computation. 
Guided by this allocation law, we scale Engram to a 27B-parameter model. Compared to a strictly iso-parameter and iso-FLOPs MoE-27B baseline, Engram-27B achieves superior efficiency across diverse domains. Crucially, the gains are not limited to knowledge-intensive tasks (e.g., MMLU: $+3.4$; CMMLU: $+4.0$; MMLU-Pro: $+1.8$), where memory capacity is intuitively beneficial; we observe even more significant improvements in general reasoning (e.g., BBH: $+5.0$; ARC-Challenge: $+3.7$; DROP: $+3.3$) and code/math domains (e.g., HumanEval: $+3.0$; MATH: $+2.4$; GSM8K: $+2.2$). 

Mechanistic analysis via LogitLens~\citep{nostalgebraist2020logitlens} and CKA~\citep{hendrycks2020measuring} reveals the source of these gains: Engram relieves the backbone from reconstructing static knowledge in early layers, thereby increasing effective depth available for complex reasoning.
Furthermore, by delegating local dependencies to lookups, Engram frees up attention capacity to focus on global context, enabling exceptional performance in long-context scenarios—substantially outperforming baselines on LongPPL~\citep{fangwrong} and RULER~\citep{hsiehruler} (e.g., Multi-Query NIAH: $97.0$ vs.\ $84.2$; Variable Tracking: $89.0$ vs.\ $77.0$).

Finally, we establish infrastructure-aware efficiency as a first-class principle. Unlike MoE's dynamic routing, Engram employs deterministic IDs to enable runtime prefetching, overlapping communication with computation. Empirical results show that offloading a 100B-parameter table to host memory incurs negligible overhead ($<3\%$). This demonstrates that Engram effectively bypasses GPU memory constraints, facilitating aggressive parameter expansion.